\documentclass[../../main.tex]{subfiles}

\begin{document}

\section{单调函数的可微性: Lebesgue 定理}\label{section:ch06-Lebesgue定理}

闭有界区间 $[c,d]$ 称为非退化的, 若 $c<d$.

\begin{definition}[Vitali 覆盖]\label{definition:ch06-Vitali覆盖}
闭有界非退化区间族 $\mathcal{F}$ 称为在 Vitali 的意义下覆盖集合 $E$, 若对 $E$ 中的每个点 $x$ 和 $\varepsilon>0$, 存在 $\mathcal{F}$ 中的区间 $I$, 它包含 $x$ 且 $\ell(I)<\varepsilon$.
\end{definition}

\begin{lemma}[Vitali 覆盖引理]\label{lemma:ch06-Vitali}
令 $E$ 为有限外测度的集合, 而 $\mathcal{F}$ 是在 Vitali 意义下覆盖 $E$ 的非退化闭有界区间族. 则对每个 $\varepsilon>0$, 存在 $\mathcal{F}$ 的有限不交子族 $\{I_{k}\}_{k=1}^{n}$ 使得
\begin{align}
m^{*}\left[E\sim\bigcup_{k=1}^{n}I_{k}\right]<\varepsilon.\label{align:ch06-Vitali}
\end{align}
\end{lemma}

\begin{proof}
由于 $m^{*}(E)<\infty$, 存在包含 $E$ 的开集 $\mathcal{O}$ 使得 $m(\mathcal{O})<\infty$. 由于 $\mathcal{F}$ 是 $E$ 的一个 Vitali 覆盖, 我们可以假设 $\mathcal{F}$ 中的每个区间包含于 $\mathcal{O}$. 根据测度的可数可加性和单调性,
\begin{align}
\text{若 }\{I_{k}\}_{k=1}^{\infty}\subseteq\mathcal{F}\text{ 是不交的},\quad\text{则}\sum_{k=1}^{\infty}\ell(I_{k})\leqslant m(\mathcal{O})<\infty.\label{align:ch06-长度和}
\end{align}
此外, 由于每个 $I_{k}$ 是闭的而 $\mathcal{F}$ 是 $E$ 的 Vitali 覆盖,
\begin{align}
\text{若 }\{I_{k}\}_{k=1}^{n}\subseteq\mathcal{F},\quad\text{则 }E\sim\bigcup_{k=1}^{n}I_{k}\subseteq\bigcup_{I\in\mathcal{F}_{n}}I,\quad\text{其中 }\mathcal{F}_{n}=\left\{I\in\mathcal{F}\mid I\cap\bigcup_{k=1}^{n}I_{k}=\varnothing\right\}.\label{align:ch06-剩余覆盖}
\end{align}
若存在 $E$ 的覆盖 $\mathcal{F}$ 的有限不交子族, 证明就完成了. 否则, 我们归纳地选取 $\mathcal{F}$ 的不交可数子族 $\{I_{k}\}_{k=1}^{\infty}$, 它具有以下性质:
\begin{align}
\text{对所有 }n,\quad E\sim\bigcup_{k=1}^{n}I_{k}\subseteq\bigcup_{k=n+1}^{\infty}5*I_{k},\label{align:ch06-5倍包含}
\end{align}
其中, 对于闭有界区间 $I$, $5*I$ 表示与 $I$ 有相同中点而长度是其 $5$ 倍的闭区间. 为开始选取, 令 $I_{1}$ 为 $\mathcal{F}$ 中的任何区间. 假定 $n$ 是自然数而 $\mathcal{F}$ 的有限不交子族 $\{I_{k}\}_{k=1}^{n}$ 已被选取. 由于 $E\sim\bigcup_{k=1}^{n}I_{k}\neq\varnothing$, \eqref{align:ch06-剩余覆盖}中定义的族 $\mathcal{F}_{n}$ 是非空的. 此外, $\mathcal{F}_{n}$ 中区间的长度的上确界 $s_{n}$ 是有限的, 这是由于 $m(\mathcal{O})$ 是这些长度的一个上界. 选取 $I_{n+1}$ 为 $\mathcal{F}_{n}$ 中使得 $\ell(I_{n+1})>s_{n}/2$ 的区间. 这归纳地定义了 $\{I_{k}\}_{k=1}^{\infty}$, 即 $\mathcal{F}$ 的一个可数不交子族使得对每个 $n$,
\begin{align}
\text{若 }I\in\mathcal{F}\text{ 且 }I\cap\bigcup_{k=1}^{n}I_{k}=\varnothing,\quad\ell(I_{n+1})>\ell(I)/2.\label{align:ch06-2倍}
\end{align}
我们从\eqref{align:ch06-长度和}推出 $\{\ell(I_{k})\}\to 0$. 固定自然数 $n$. 为证明包含关系\eqref{align:ch06-5倍包含}, 令 $x$ 属于 $E\sim\bigcup_{k=1}^{n}I_{k}$. 我们从\eqref{align:ch06-剩余覆盖}推出存在 $I\in\mathcal{F}$, 它包含 $x$ 且与 $\bigcup_{k=1}^{n}I_{k}$ 不交. 现在 $I$ 必须与某个 $I_{k}$ 有非空交, 否则根据\eqref{align:ch06-2倍}, 对所有 $k$, $\ell(I_{k})>\ell(I)/2$, 这与 $\{\ell(I_{k})\}$ 收敛到 $0$ 矛盾. 令 $N$ 为使得 $I\cap I_{N}\neq\varnothing$ 的第一个自然数. 则 $N>n$. 由于 $I\cap\bigcup_{k=1}^{N-1}I_{k}=\varnothing$, 我们从\eqref{align:ch06-2倍}推出 $\ell(I_{N})>\ell(I)/2$. 由于 $x$ 属于 $I$ 且 $I\cap I_{N}\neq\varnothing$, $x$ 到 $I_{N}$ 的中点的距离至多为 $\ell(I)+1/2\cdot\ell(I_{N})$. 因此, 由于 $\ell(I)<2\cdot\ell(I_{N})$, $x$ 到 $I_{N}$ 的中点的距离小于 $5/2\cdot\ell(I_{N})$. 这意味着 $x$ 属于 $5*I_{N}$. 因此,
\begin{align*}
x\in 5*I_{N}\subseteq\bigcup_{k=n+1}^{\infty}5*I_{k}.
\end{align*}
我们证明了包含关系\eqref{align:ch06-5倍包含}.

令 $\varepsilon>0$. 我们从\eqref{align:ch06-长度和}推出存在自然数 $n$ 使得 $\sum_{k=n+1}^{\infty}\ell(I_{k})<\varepsilon/5$. $n$ 的这个选择与包含关系\eqref{align:ch06-5倍包含}, 以及测度的单调性、可数可加性一起, 证明了\eqref{align:ch06-Vitali}.
\end{proof}

对于实值函数 $f$ 和它的定义域的一个内点 $x$, $f$ 在 $x$ 的上导数 $\overline{D}f(x)$ 与下导数 $\underline{D}f(x)$ 定义如下:
\begin{align*}
\overline{D}f(x)=\lim_{h\to 0}\left[\sup_{0<|t|\leqslant h}\frac{f(x+t)-f(x)}{t}\right]
\end{align*}
与
\begin{align*}
\underline{D}f(x)=\lim_{h\to 0}\left[\inf_{0<|t|\leqslant h}\frac{f(x+t)-f(x)}{t}\right].
\end{align*}
我们有 $\overline{D}f(x)\geqslant\underline{D}f(x)$. 若 $\overline{D}f(x)$ 等于 $\underline{D}f(x)$ 且是有限的, 我们说 $f$ 在 $x$ 是可微的且定义 $f'(x)$ 为上导数与下导数的公共值.

微积分的中值定理告诉我们, 若函数 $f$ 在闭有界区间 $[c,d]$ 上连续且在它的内部 $(c,d)$ 可微, 满足在 $(c,d)$ 上 $f'\geqslant\alpha$, 则
\begin{align*}
\alpha\cdot(d-c)\leqslant[f(d)-f(c)].
\end{align*}
该不等式的以下推广的证明很好地说明了 Vitali 覆盖引理与函数的单调性质的富有成效的交互.

\begin{lemma}\label{lemma:ch06-上导数估计}
令 $f$ 为闭有界区间 $[a,b]$ 上的递增函数. 则对每个 $\alpha>0$,
\begin{align}
m^{*}\{x\in(a,b)\mid\overline{D}f(x)\geqslant\alpha\}\leqslant\frac{1}{\alpha}\cdot[f(b)-f(a)]\label{align:ch06-上导数估计}
\end{align}
且
\begin{align}
m^{*}\{x\in(a,b)\mid\overline{D}f(x)=\infty\}=0.\label{align:ch06-上导数无穷}
\end{align}
\end{lemma}

\begin{proof}
令 $\alpha>0$. 定义 $E_{\alpha}=\{x\in(a,b)\mid\overline{D}f(x)\geqslant\alpha\}$. 选取 $\alpha'\in(0,\alpha)$. 令 $\mathcal{F}$ 为包含于 $(a,b)$ 的闭有界区间 $[c,d]$ 所组成的族, 使得 $f(d)-f(c)\geqslant\alpha'(d-c)$. 由于在 $E_{\alpha}$ 上 $\overline{D}f\geqslant\alpha$, $\mathcal{F}$ 是 $E_{\alpha}$ 的一个 Vitali 覆盖. Vitali 覆盖引理告诉我们存在 $\mathcal{F}$ 的有限不交子族 $\{[c_{k},d_{k}]\}_{k=1}^{n}$ 使得
\begin{align*}
m^{*}\left[E_{\alpha}\sim\bigcup_{k=1}^{n}[c_{k},d_{k}]\right]<\varepsilon.
\end{align*}
由于 $E_{\alpha}\subseteq\bigcup_{k=1}^{n}[c_{k},d_{k}]\cup\{E_{\alpha}\sim\bigcup_{k=1}^{n}[c_{k},d_{k}]\}$, 根据外测度的有限次可加性, 以及前一个不等式和区间 $[c_{k},d_{k}]$ 的选取,
\begin{align}
m^{*}(E_{\alpha})\leqslant\sum_{k=1}^{n}(d_{k}-c_{k})+\varepsilon\leqslant\frac{1}{\alpha'}\cdot\sum_{k=1}^{n}[f(d_{k})-f(c_{k})]+\varepsilon.\label{align:ch06-估计1}
\end{align}
然而, 函数 $f$ 在 $[a,b]$ 上是递增的, 而 $\{[c_{k},d_{k}]\}_{k=1}^{n}$ 是 $[a,b]$ 的不交子区间族. 于是,
\begin{align*}
\sum_{k=1}^{n}[f(d_{k})-f(c_{k})]\leqslant f(b)-f(a).
\end{align*}
因此对每个 $\varepsilon>0$ 和每个 $\alpha'\in(0,\alpha)$,
\begin{align*}
m^{*}(E_{\alpha})\leqslant\frac{1}{\alpha'}\cdot[f(b)-f(a)]+\varepsilon.
\end{align*}
这证明了\eqref{align:ch06-上导数估计}. 对每个自然数 $n$, $\{x\in(a,b)\mid\overline{D}f(x)=\infty\}\subseteq E_{n}$, 因此
\begin{align*}
m^{*}\{x\in(a,b)\mid\overline{D}f(x)=\infty\}\leqslant m^{*}(E_{n})\leqslant\frac{1}{n}\cdot(f(b)-f(a)).
\end{align*}
这证明了\eqref{align:ch06-上导数无穷}.
\end{proof}

\begin{theorem}[Lebesgue 定理]\label{theorem:ch06-单调可微}
若函数 $f$ 在开区间 $(a,b)$ 上单调, 则它在 $(a,b)$ 上是几乎处处可微的.
\end{theorem}

\begin{proof}
假设 $f$ 是递增的. 此外, 假设 $(a,b)$ 是有界的. 否则, 将 $(a,b)$ 表示为开有界区间的上升序列的并集, 进而用 Lebesgue 测度的连续性. $(a,b)$ 中那些使得 $\overline{D}f(x)>\underline{D}f(x)$ 的点 $x$ 的集合是集合
\begin{align*}
E_{\alpha,\beta}=\{x\in(a,b)\mid\overline{D}f(x)>\alpha>\beta>\underline{D}f(x)\}
\end{align*}
的并集, 其中 $\alpha$ 和 $\beta$ 是有理数. 因此, 由于这是可数族, 根据外测度的可数次可加性, 仅须证明每个 $E_{\alpha,\beta}$ 的外测度为零. 固定有理数 $\alpha$ 和 $\beta$ 使得 $\alpha>\beta$, 且设 $E=E_{\alpha,\beta}$. 令 $\varepsilon>0$. 选取开集 $\mathcal{O}$ 使得
\begin{align}
E\subseteq\mathcal{O}\subseteq(a,b)\quad\text{且}\quad m(\mathcal{O})<m^{*}(E)+\varepsilon.\label{align:ch06-开集}
\end{align}
令 $\mathcal{F}$ 为包含于 $\mathcal{O}$ 的闭有界区间 $[c,d]$ 所组成的族, 使得 $f(d)-f(c)<\beta(d-c)$. 由于在 $E$ 上 $\underline{D}f<\beta$, $\mathcal{F}$ 是 $E$ 的一个 Vitali 覆盖. Vitali 覆盖引理告诉我们存在 $\mathcal{F}$ 的有限不交子族 $\{[c_{k},d_{k}]\}_{k=1}^{n}$ 使得
\begin{align}
m^{*}\left[E\sim\bigcup_{k=1}^{n}[c_{k},d_{k}]\right]<\varepsilon.\label{align:ch06-剩余小}
\end{align}
根据区间 $[c_{k},d_{k}]$ 的选取、$\mathcal{O}$ 中的不交区间族 $\{[c_{k},d_{k}]\}_{k=1}^{n}$ 的并的包含关系以及\eqref{align:ch06-开集},
\begin{align}
\sum_{k=1}^{n}[f(d_{k})-f(c_{k})]<\beta\left[\sum_{k=1}^{n}(d_{k}-c_{k})\right]\leqslant\beta\cdot m(\mathcal{O})\leqslant\beta\cdot[m^{*}(E)+\varepsilon].\label{align:ch06-导数和}
\end{align}
对于 $1\leqslant k\leqslant n$, 我们从前一个引理应用于 $f$ 在 $[c_{k},d_{k}]$ 的限制推出
\begin{align*}
m^{*}(E\cap(c_{k},d_{k}))\leqslant\frac{1}{\alpha}[f(d_{k})-f(c_{k})].
\end{align*}
因此, 根据\eqref{align:ch06-剩余小},
\begin{align}
m^{*}(E)\leqslant\sum_{k=1}^{n}m^{*}(E\cap(c_{k},d_{k}))+\varepsilon\leqslant\frac{1}{\alpha}\left[\sum_{k=1}^{n}[f(d_{k})-f(c_{k})]\right]+\varepsilon.\label{align:ch06-估计2}
\end{align}
我们从\eqref{align:ch06-导数和}和\eqref{align:ch06-估计2}推出, 对所有 $\varepsilon>0$,
\begin{align*}
m^{*}(E)\leqslant\frac{\beta}{\alpha}\cdot m^{*}(E)+\frac{1}{\alpha}\cdot\varepsilon+\varepsilon.
\end{align*}
因此, 由于 $0\leqslant m^{*}(E)<\infty$ 且 $\beta/\alpha<1$, $m^{*}(E)=0$.
\end{proof}

Lebesgue 定理在以下意义是最佳的, 若 $E$ 是包含于开区间 $(a,b)$ 的测度为零的集合, 存在 $(a,b)$ 上的递增函数, 它在 $E$ 中的每个点处不是可微的 (见习题 10).

令 $f$ 在闭有界区间 $[a,b]$ 上可积. 延拓 $f$ 使得它在 $(b,b+1)$ 上取值 $f(b)$. 对 $0<h\leqslant 1$, 定义 $[a,b]$ 的均差函数 $\operatorname{Diff}_{h}f$ 与平均值函数 $\operatorname{Avg}_{h}f$ 为
\begin{align*}
\text{对所有 }x\in[a,b],\quad\operatorname{Diff}_{h}f(x)=\frac{f(x+h)-f(x)}{h}\quad\text{与}\quad\operatorname{Avg}_{h}f(x)=\frac{1}{h}\int_{x}^{x+h}f.
\end{align*}
在积分中换元和抵消掉一些项, 对所有 $a\leqslant u<v\leqslant b$,
\begin{align}
\int_{u}^{v}\operatorname{Diff}_{h}f=\operatorname{Avg}_{h}f(v)-\operatorname{Avg}_{h}f(u).\label{align:ch06-换元}
\end{align}

\begin{corollary}\label{corollary:ch06-单调导数可积}
令 $f$ 为闭有界区间 $[a,b]$ 上的递增函数. 则 $f'$ 在 $[a,b]$ 上可积且
\begin{align}
\int_{a}^{b}f'\leqslant f(b)-f(a).\label{align:ch06-导数积分}
\end{align}
\end{corollary}

\begin{proof}
由于 $f$ 在 $[a,b+1]$ 上递增, 它是可测的 (见习题 22), 因此均差函数也是可测的. Lebesgue 定理告诉我们 $f$ 在 $(a,b)$ 上几乎处处可微. 因此 $\{\operatorname{Diff}_{1/n}f\}$ 是在 $[a,b]$ 上几乎处处逐点收敛于 $f'$ 的非负可测函数列. 根据 Fatou 引理,
\begin{align}
\int_{a}^{b}f'\leqslant\liminf_{n\to\infty}\left[\int_{a}^{b}\operatorname{Diff}_{1/n}f\right].\label{align:ch06-Fatou}
\end{align}
根据换元公式\eqref{align:ch06-换元}, 对每个自然数 $n$, 由于 $f$ 是递增的,
\begin{align*}
\int_{a}^{h}\operatorname{Diff}_{1/n}f=\frac{1}{1/n}\cdot\int_{b}^{b+1/n}f-\frac{1}{1/n}\cdot\int_{a}^{a+1/n}f=f(b)-\frac{1}{1/n}\cdot\int_{a}^{a+1/n}f\leqslant f(b)-f(a).
\end{align*}
因此
\begin{align}
\limsup_{n\to\infty}\left[\int_{a}^{b}\operatorname{Diff}_{1/n}f\right]\leqslant f(b)-f(a).\label{align:ch06-上界}
\end{align}
不等式\eqref{align:ch06-导数积分}从不等式\eqref{align:ch06-Fatou}和\eqref{align:ch06-上界}得出.
\end{proof}

不等式\eqref{align:ch06-导数积分}中的积分与 $f$ 在端点的取值无关. 另一方面, 该等式的右边对开有界区间 $(a,b)$ 上 $f$ 到其定义域的闭包 $[a,b]$ 的任何递增延拓成立. 因此等式\eqref{align:ch06-导数积分}的一个更紧凑形式是
\begin{align}
\int_{a}^{b}f'\leqslant\sup_{x\in(a,b)}f(x)-\inf_{x\in(a,b)}f(x).\label{align:ch06-紧凑不等式}
\end{align}

不等式的右边等于 $f(b)-f(a)$ 当且仅当 $f$ 在端点连续. 然而, 即使 $f$ 在 $[a,b]$ 上是递增的且连续, 不等式\eqref{align:ch06-导数积分}也可以是严格的. 对于 $[0,1]$ 上的 Cantor-Lebesgue 函数 $\varphi$ 它是严格的, 这是由于 $\varphi(1)-\varphi(0)=1$ 而 $\varphi'$ 在 $(0,1)$ 上几乎处处为零. 我们证明对任何 $[a,b]$ 上的递增函数 $f$, \eqref{align:ch06-导数积分}是等式当且仅当函数在 $[a,b]$ 上是绝对连续的 (见即将介绍的推论\cref{corollary:ch06-单调绝对连续}).

\begin{exercise}\label{exercise:ch06-05}
证明 Vitali 覆盖引理不能推广到覆盖族有退化闭区间的情形.
\end{exercise}

\begin{exercise}\label{exercise:ch06-06}
证明 Vitali 覆盖引理能推广到覆盖族由非退化的一般区间组成的情形.
\end{exercise}

\begin{exercise}\label{exercise:ch06-07}
令 $f$ 在 $\mathbf{R}$ 上连续. 是否有开区间使得 $f$ 在其上是单调的?
\end{exercise}

\begin{exercise}\label{exercise:ch06-08}
令 $I$ 和 $J$ 为闭有界区间而 $\gamma>0$ 使得 $\ell(I)>\gamma\ell(J)$. 假设 $I\cap J\neq\varnothing$. 证明: 若 $\gamma\geqslant 1/2$, 则 $J\subseteq 5*I$, 其中 $5*I$ 表示与 $I$ 有相同的中心而长度是其 $5$ 倍的区间. 若 $0<\gamma<1/2$, 同样的结论成立吗?
\end{exercise}

\begin{exercise}\label{exercise:ch06-09}
证明实数集 $E$ 的测度为零当且仅当存在开区间的可数族 $\{I_{k}\}_{k=1}^{\infty}$, 使得 $E$ 中的每个点属于无穷多个区间 $I_{k}$ 且 $\sum_{k=1}^{\infty}\ell(I_{k})<\infty$.
\end{exercise}

\begin{exercise}\label{exercise:ch06-10}
(Riesz-Nagy) 令 $E$ 为包含于开区间 $(a,b)$ 的测度为零的集合. 根据前一个习题, 存在包含于 $(a,b)$ 的开区间的可数族 $\{(c_{k},d_{k})\}_{k=1}^{\infty}$, 使得 $E$ 中的每个点属于该族中的无穷多个区间且 $\sum_{k=1}^{\infty}(d_{k}-c_{k})<\infty$. 对所有 $(a,b)$ 中的 $x$ 定义
\begin{align*}
f(x)=\sum_{k=1}^{\infty}\ell((c_{k},d_{k})\cap(-\infty,x)).
\end{align*}
证明 $f$ 是递增的且在 $E$ 中的每个点处不可微.
\end{exercise}

\begin{exercise}\label{exercise:ch06-11}
对于实数 $\alpha<\beta$ 和 $\gamma>0$, 证明: 若 $g$ 在 $[\alpha+\gamma,\beta+\gamma]$ 上可积, 则
\begin{align*}
\int_{\alpha}^{\beta}g(t+\gamma)\,dt=\int_{\alpha+\gamma}^{\beta+\gamma}g(t)\,dt.
\end{align*}
通过相继考虑简单函数、有界可测函数、非负可积函数、一般的可积函数证明该换元公式. 用它证明\eqref{align:ch06-换元}.
\end{exercise}

\begin{exercise}\label{exercise:ch06-12}
计算有理数集的特征函数的上导数与下导数.
\end{exercise}

\begin{exercise}\label{exercise:ch06-13}
令 $E$ 为有限外测度的集合而 $\mathcal{F}$ 是 Vitali 意义下覆盖 $E$ 的闭有界区间族. 证明存在 $\mathcal{F}$ 中的区间的可数不交族 $\{I_{k}\}_{k=1}^{\infty}$ 使得
\begin{align*}
m^{*}\left[E\sim\bigcup_{k=1}^{\infty}I_{k}\right]=0.
\end{align*}
\end{exercise}

\begin{exercise}\label{exercise:ch06-14}
用 Vitali 覆盖引理证明任何闭有界非退化区间族 (可数或不可数) 的并是可测的.
\end{exercise}

\begin{exercise}\label{exercise:ch06-15}
定义 $\mathbf{R}$ 上的 $f$ 为
\begin{align*}
f(x)=
\begin{cases}
x\sin(1/x)&\text{若 }x\neq 0\\
0&\text{若 }x=0
\end{cases}
\end{align*}
找出 $f$ 在 $x=0$ 的上导数与下导数.
\end{exercise}

\begin{exercise}\label{exercise:ch06-16}
令 $g$ 在 $[a,b]$ 上可积. 定义 $g$ 的反导数为在 $[a,b]$ 上定义的函数 $f$,
\begin{align*}
f(x)=\int_{a}^{x}g,\quad\text{对所有 }x\in[a,b].
\end{align*}
证明 $f$ 在 $(a,b)$ 上几乎处处可微.
\end{exercise}

\begin{exercise}\label{exercise:ch06-17}
令 $f$ 为开有界区间 $(a,b)$ 上递增的有界函数. 证明\eqref{align:ch06-紧凑不等式}.
\end{exercise}

\begin{exercise}\label{exercise:ch06-18}
证明: 若 $f$ 定义在 $(a,b)$ 上而 $c\in(a,b)$ 是 $f$ 的局部极小点, 则 $\underline{D}f(c)\leqslant 0\leqslant\overline{D}f(c)$.
\end{exercise}

\begin{exercise}\label{exercise:ch06-19}
令 $f$ 在 $[a,b]$ 上连续, 满足在 $(a,b)$ 上 $\overline{D}f\geqslant 0$. 证明 $f$ 在 $[a,b]$ 上递增. (提示: 首先对 $(a,b)$ 上满足 $\underline{D}g\geqslant\varepsilon>0$ 的函数 $g$ 证明该结论. 将此用于函数 $g(x)=f(x)+\varepsilon x$.)
\end{exercise}

\begin{exercise}\label{exercise:ch06-20}
令 $f$ 和 $g$ 为 $(a,b)$ 上的实值函数. 证明在 $(a,b)$ 上,
\begin{align*}
\underline{D}f+\underline{D}g\leqslant\underline{D}(f+g)\leqslant\overline{D}(f+g)\leqslant\overline{D}f+\overline{D}g.
\end{align*}
\end{exercise}

\begin{exercise}\label{exercise:ch06-21}
令 $f$ 定义在 $[a,b]$ 上而 $g$ 是 $[a,\beta]$ 上的连续函数, 它在 $\gamma\in(a,\beta)$ 可微, 其中 $g(\gamma)=c\in(a,b)$. 证明以下命题.
\begin{enumerate}[(i)]
\item 若 $g'(\gamma)>0$, 则 $\overline{D}(f\circ g)(\gamma)=\overline{D}f(c)g'(\gamma)$.
\item 若 $g'(\gamma)=0$ 且 $f$ 在 $c$ 的上导数与下导数是有限的, 则 $\overline{D}(f\circ g)(\gamma)=0$.
\end{enumerate}
\end{exercise}

\begin{exercise}\label{exercise:ch06-22}
证明定义在区间上的严格递增函数是可测的, 并用此证明定义在区间上的单调函数是可测的.
\end{exercise}

\begin{exercise}\label{exercise:ch06-23}
证明 $[a,b]$ 上的连续函数 $f$ 是 Lipschitz 的, 若它的下导数与上导数在 $(a,b)$ 上有界.
\end{exercise}

\begin{exercise}\label{exercise:ch06-24}
证明对于本节最后的注中定义的 $f$, $f'$ 在 $[0,1]$ 上不可积.
\end{exercise}

\end{document}